\section{Challenge 3: Nhân hai ma trận}
\label{sec:challenge3}

\begin{requirementbox}{Yêu cầu bài toán}
Cho hai ma trận vuông $A$ và $B$ kích thước $N \times N$, tính $C = A \times B$, đồng thời trả về đường chéo chính, đường chéo phụ và tổng toàn bộ phần tử của $C$.
\end{requirementbox}

\subsection{Cách hiện thực trong mã nguồn}
File \codepath{challenge/challenge3.py} đọc số nguyên bằng generator \func{\_int\_stream} để giảm chi phí split toàn bộ chuỗi. Ma trận $B$ được lưu ở dạng chuyển vị \func{\_BT\_DATA}; nhờ đó tích vô hướng giữa một dòng của $A$ và một cột của $B$ trở thành truy cập tuần tự giữa hai list. Hàm \func{\_mul\_chunk(start\_row, end\_row)} tính các dòng $C[start\_row:end\_row]$, đồng thời trả về đoạn đường chéo chính, đường chéo phụ và tổng cục bộ.

\subsection{Phân tích song song hóa}
Mỗi dòng của ma trận kết quả $C$ có thể tính độc lập sau khi đã có $A$ và $B^T$. Vì vậy chiến lược chia việc tự nhiên là chia theo cụm dòng:
\begin{itemize}
  \item Worker nhận một khoảng dòng $[s, e)$.
  \item Worker đọc toàn bộ $A$ và $B^T$ nhưng chỉ ghi/trả về các dòng $C$ thuộc khoảng của mình.
  \item Tiến trình chính ghép các dòng theo \func{start\_row}, ghép đường chéo và cộng tổng cục bộ.
\end{itemize}

\subsection{Thiết kế thuật toán song song}
\begin{lstlisting}[caption={Pseudocode song song cho Challenge 3}]
MAIN(path):
    N = read_N(path)
    A = read_matrix_A(path, N)
    BT = read_transposed_matrix_B(path, N)
    tasks = split_rows_into_chunks(N, workers, chunks_per_worker=4)

    for each (start_row, end_row) in parallel:
        rows = []
        diag_main_part = []
        diag_secondary_part = []
        local_sum = 0

        for i in range(start_row, end_row):
            for j in range(N):
                C[i][j] = dot(A[i], BT[j])
            update diagonals and local_sum

    merge rows, diagonals, and sums
    return {matrix, diag_main, diag_secondary, total_sum}
\end{lstlisting}

\subsection{Dữ liệu, phụ thuộc và chi phí}
\begin{itemize}
  \item \textbf{Phân phối dữ liệu}: $A$ và $B^T$ là dữ liệu đọc chung qua biến toàn cục trong worker; kết quả được trả về theo chunk.
  \item \textbf{Phụ thuộc dữ liệu}: các dòng $C_i$ độc lập; không có worker nào cần kết quả của worker khác.
  \item \textbf{Cân bằng tải}: chia theo dòng cho tải gần đều vì mỗi dòng cần $N^2$ phép nhân-cộng.
  \item \textbf{Chi phí giao tiếp}: cao hơn các challenge khác vì worker trả về nhiều dòng của ma trận $C$; với $N$ lớn, chi phí bộ nhớ và pickle/IPC có thể chi phối.
  \item \textbf{Giới hạn bộ nhớ}: việc trả về toàn bộ $C$ có độ phức tạp bộ nhớ \bigO{N^2}; khi thử $N=5000$ cần ghi rõ cấu hình RAM và khả năng chạy thực tế.
\end{itemize}

\subsection{Đánh giá hiệu năng}
\begin{table}[H]
\centering
\caption{Kết quả đo hiệu năng cho Challenge 3}
\begin{tabular}{@{}rrrrrr@{}}
\toprule
$N$ & Số worker & $T_{\mathrm{seq}}$ (s) & $T_{\mathrm{par}}$ (s) & Speedup & Efficiency \\
\midrule
128 & \placeholder{p} & \placeholder{tseq} & \placeholder{tpar} & \placeholder{speedup} & \placeholder{eff} \\
256 & \placeholder{p} & \placeholder{tseq} & \placeholder{tpar} & \placeholder{speedup} & \placeholder{eff} \\
512 & \placeholder{p} & \placeholder{tseq} & \placeholder{tpar} & \placeholder{speedup} & \placeholder{eff} \\
1024 & \placeholder{p} & \placeholder{tseq} & \placeholder{tpar} & \placeholder{speedup} & \placeholder{eff} \\
5000 & \placeholder{p} & \placeholder{tseq} & \placeholder{tpar} & \placeholder{speedup} & \placeholder{eff} \\
\bottomrule
\end{tabular}
\end{table}

\paragraph{Nhận xét.}
\placeholder{Giải thích xu hướng khi N tăng: lượng tính toán tăng nhanh nên song song hóa theo dòng có lợi hơn, nhưng chi phí bộ nhớ và trả ma trận kết quả có thể giới hạn speedup.}
